% ============================================================================
% 03_modelos_casos_uso.tex — Cuerpo del entregable "Modelos de Casos de Uso"
% Formato "fully dressed" (Larman, UML y Patrones) para cada caso de uso.
% ============================================================================
\section{Modelos de Casos de Uso}

\subsection{Introducción}

Esta sección desarrolla, en formato completamente vestido (\textit{fully
dressed}) según Craig Larman, cada uno de los once casos de uso
identificados en el Entregable de Diagrama de Casos de Uso. Los escenarios
se redactan a nivel esencial (intención del actor), sin describir botones,
menús ni otros detalles de la interfaz concreta, que se documentan en el
entregable de mockups. CU-09, CU-10 y CU-11 corresponden a los requisitos
oficiales RF-03, RF-06 y RF-04 respectivamente, y están marcados como
\textbf{diseñado --- pendiente de implementación}: no existe código que los
materialice, pero se especifican con el mismo nivel de detalle que el
resto para que el diseño quede completo.

\subsection{Actores}

\begin{table}[H]
    \centering
    \caption{Actores de LSA Salud}
    \label{tab:actores-mcu}
    \begin{tabularx}{\textwidth}{l l X}
        \toprule
        \textbf{ID} & \textbf{Actor} & \textbf{Descripción} \\
        \midrule
        ACT-01 & Paciente & Persona sorda o hipoacúsica que se comunica mediante LSA. \\
        ACT-02 & Personal de Recepción o Enfermería & Personal de recepción o enfermería que responde al paciente. \\
        ACT-03 & Desarrollador de IA / Equipo de Proyecto & Personal técnico que amplía el vocabulario y reentrena el modelo. \\
        \bottomrule
    \end{tabularx}
\end{table}

% ============================================================================
\subsection{CU-01: Comunicarse mediante Lengua de Señas Argentina}

\begin{longtable}{@{}>{\raggedright\arraybackslash}p{4.1cm} >{\raggedright\arraybackslash}p{\dimexpr\textwidth-4.1cm-2\tabcolsep\relax}@{}}
\caption{CU-01 --- Comunicarse mediante Lengua de Señas Argentina}\\
\toprule
\textbf{Nombre} & Comunicarse mediante Lengua de Señas Argentina \\
\textbf{ID} & CU-01 \\
\textbf{Actor principal} & Paciente \\
\textbf{Personal involucrado e intereses} &
Paciente: quiere expresar una necesidad médica sin depender de un
intérprete humano. \newline
Personal de Recepción o Enfermería: necesita recibir el mensaje del
paciente en español, de forma clara y oportuna. \\
\textbf{Precondiciones} &
El paciente otorgó el consentimiento explícito de cámara (CU-09).
Estación del Paciente abierta, cámara autorizada, backend en ejecución,
modelo de reconocimiento cargado. \\
\textbf{Garantías de éxito} &
El mensaje traducido queda registrado en la conversación compartida y, si
corresponde, se reprodujo en voz. \\
\textbf{Escenario principal de éxito} &
\begin{enumerate}[leftmargin=*, itemsep=2pt]
    \item El paciente se ubica frente a la cámara y comienza a hacer señas.
    \item El sistema captura de forma continua los movimientos del paciente.
    \item El sistema reconoce cada seña realizada y la agrega a la
    secuencia en curso cuando la reconoce con suficiente certeza.
    \item El paciente hace una pausa (baja las manos) al terminar de
    expresar su idea.
    \item El sistema detecta la pausa y da por cerrada la secuencia de
    señas.
    \item El sistema traduce la secuencia a una oración en español natural
    (CU-02: Traducir Secuencia de Señas a Oración Natural).
    \item El sistema sintetiza la oración en voz (CU-03: Sintetizar
    Oración en Voz), si la síntesis está disponible.
    \item El sistema agrega el mensaje traducido a la conversación,
    visible para el personal de recepción o enfermería.
    \item El sistema registra un evento de auditoría anónimo de la
    operación (CU-11: Registrar Evento de Auditoría).
\end{enumerate} \\
\textbf{Extensiones} &
\begin{itemize}[leftmargin=*, itemsep=2pt]
    \item[3a.] La seña reconocida no alcanza el nivel de certeza mínimo: el
    sistema la descarta y continúa observando.
    \item[3b.] La seña reconocida es el estado de reposo (``idle''): el
    sistema no la agrega a la secuencia.
    \item[6a.] El sistema no logra traducir la secuencia en una oración
    clara: descarta el resultado y no genera mensaje.
    \item[7a.] La síntesis de voz no está disponible: el sistema continúa
    sin reproducir audio.
\end{itemize} \\
\textbf{Requisitos especiales} & RNF-T-01, RNF-T-05, RNF-01.1 (ver ERS, sección 3.5). \\
\textbf{Variantes de tecnología y datos} &
Captura de landmarks por MediaPipe Holistic (por defecto) o MediaPipe
Tasks (Pose + Mano + Cara, activable con el parámetro \texttt{mp=tasks});
transporte de los frames por WebSocket en formato JSON (v1) o binario
(v2). \\
\textbf{Frecuencia} &
Continua durante toda la sesión de atención; se repite cada vez que el
paciente forma una idea completa. \\
\textbf{Temas abiertos} &
Vocabulario limitado a las clases entrenadas; sin soporte para
reconocimiento continuo de LSA como lengua completa (fuera de alcance
actual). \\
\bottomrule
\end{longtable}

% ============================================================================
\subsection{CU-02: Traducir Secuencia de Señas a Oración Natural}

\begin{longtable}{@{}>{\raggedright\arraybackslash}p{4.1cm} >{\raggedright\arraybackslash}p{\dimexpr\textwidth-4.1cm-2\tabcolsep\relax}@{}}
\caption{CU-02 --- Traducir Secuencia de Señas a Oración Natural}\\
\toprule
\textbf{Nombre} & Traducir Secuencia de Señas a Oración Natural \\
\textbf{ID} & CU-02 \\
\textbf{Actor principal} & Paciente (caso de uso incluido desde CU-01) \\
\textbf{Personal involucrado e intereses} &
Paciente: quiere que su idea se entienda correctamente, sin agregados
inventados. \newline
Personal de Recepción o Enfermería: necesita una oración clara y fiel a lo
expresado. \\
\textbf{Precondiciones} &
Existe una secuencia de señas cerrada, con al menos una palabra
reconocida. \\
\textbf{Garantías de éxito} &
Se obtuvo una oración en español, o se determinó que la secuencia no
forma una idea clara. \\
\textbf{Escenario principal de éxito} &
\begin{enumerate}[leftmargin=*, itemsep=2pt]
    \item El sistema recibe la secuencia de palabras sueltas.
    \item El sistema normaliza la secuencia (descarta repeticiones y la
    palabra de reposo).
    \item El sistema busca si la combinación coincide con una plantilla
    médica conocida.
    \item El sistema toma la frase de la plantilla como traducción.
    \item El sistema registra la traducción obtenida, para reutilizarla si
    la misma combinación vuelve a aparecer.
\end{enumerate} \\
\textbf{Extensiones} &
\begin{itemize}[leftmargin=*, itemsep=2pt]
    \item[3a.] No hay coincidencia con ninguna plantilla: el sistema busca
    la combinación entre las traducciones generadas recientemente
    (3a.1); si tampoco encuentra coincidencia, delega la construcción de
    la oración a un modelo de lenguaje local, indicándole que respete la
    persona gramatical y que no agregue síntomas no mencionados (3a.2).
    \item[3a.2a.] El modelo de lenguaje no responde a tiempo o falla: el
    sistema usa como traducción las palabras originales, sin interrumpir
    el flujo.
    \item[3a.2b.] El modelo de lenguaje determina que la combinación no
    forma una idea clara: el sistema descarta la traducción y no genera
    mensaje.
\end{itemize} \\
\textbf{Requisitos especiales} & RNF-T-01, RF-01.5. \\
\textbf{Variantes de tecnología y datos} &
Plantillas médicas precalculadas (diccionario en el backend); caché de
traducciones recientes; modelo de lenguaje local Qwen2 1.5B servido por
Ollama en \texttt{localhost:11434}. \\
\textbf{Frecuencia} & Una vez por cada secuencia de señas cerrada. \\
\textbf{Temas abiertos} &
Cobertura de plantillas limitada a las clases entrenadas actualmente; la
restricción de no inventar síntomas se indica en el prompt del modelo de
lenguaje, pero no se verifica automáticamente sobre su respuesta. \\
\bottomrule
\end{longtable}

% ============================================================================
\subsection{CU-03: Sintetizar Oración en Voz}

\begin{longtable}{@{}>{\raggedright\arraybackslash}p{4.1cm} >{\raggedright\arraybackslash}p{\dimexpr\textwidth-4.1cm-2\tabcolsep\relax}@{}}
\caption{CU-03 --- Sintetizar Oración en Voz}\\
\toprule
\textbf{Nombre} & Sintetizar Oración en Voz \\
\textbf{ID} & CU-03 \\
\textbf{Actor principal} & Paciente (caso de uso incluido desde CU-01) \\
\textbf{Personal involucrado e intereses} &
Personal de Recepción o Enfermería: quiere poder escuchar el mensaje, no
solo leerlo, para atender más rápido. \\
\textbf{Precondiciones} &
Existe una oración traducida; el motor de síntesis de voz está cargado. \\
\textbf{Garantías de éxito} &
Se generó el audio de la oración, o se determinó que la síntesis no está
disponible. \\
\textbf{Escenario principal de éxito} &
\begin{enumerate}[leftmargin=*, itemsep=2pt]
    \item El sistema recibe la oración traducida.
    \item El sistema verifica si ya generó audio antes para esa misma
    oración.
    \item El sistema sintetiza el audio con voz en español rioplatense.
    \item El sistema entrega el audio para su reproducción, tras una
    breve ventana de exhibición del texto en pantalla.
\end{enumerate} \\
\textbf{Extensiones} &
\begin{itemize}[leftmargin=*, itemsep=2pt]
    \item[2a.] El audio ya fue generado antes para la misma oración: el
    sistema reutiliza el audio ya generado en lugar de sintetizarlo de
    nuevo.
    \item[3a.] El motor de síntesis de voz no está cargado: el sistema
    entrega el mensaje sin audio.
\end{itemize} \\
\textbf{Requisitos especiales} & RNF-T-01, RF-01.6. \\
\textbf{Variantes de tecnología y datos} &
Motor Piper TTS con modelo de voz en español argentino
(\texttt{es\_AR-daniela-high}); caché de audio en memoria (hasta 50
oraciones). \\
\textbf{Frecuencia} & Una vez por cada oración traducida exitosamente. \\
\textbf{Temas abiertos} & Ninguno relevado. \\
\bottomrule
\end{longtable}

% ============================================================================
\subsection{CU-04: Responder al Paciente}

\begin{longtable}{@{}>{\raggedright\arraybackslash}p{4.1cm} >{\raggedright\arraybackslash}p{\dimexpr\textwidth-4.1cm-2\tabcolsep\relax}@{}}
\caption{CU-04 --- Responder al Paciente}\\
\toprule
\textbf{Nombre} & Responder al Paciente \\
\textbf{ID} & CU-04 \\
\textbf{Actor principal} & Personal de Recepción o Enfermería \\
\textbf{Personal involucrado e intereses} &
Personal de Recepción o Enfermería: quiere responder rápido, con el método
que le resulte más cómodo. \newline
Paciente: necesita leer la respuesta en su propia pantalla sin demora. \\
\textbf{Precondiciones} &
Estación de Recepción abierta; existe una conversación activa (puede
estar vacía). \\
\textbf{Garantías de éxito} &
La respuesta del personal de recepción o enfermería queda registrada en la
conversación y visible en la Estación del Paciente. \\
\textbf{Escenario principal de éxito} &
\begin{enumerate}[leftmargin=*, itemsep=2pt]
    \item El personal redacta una respuesta en su propio idioma.
    \item El sistema valida que la respuesta no esté vacía.
    \item El personal confirma el envío.
    \item El sistema agrega la respuesta a la conversación compartida.
    \item El sistema muestra la respuesta en ambas estaciones
    sincronizadas.
    \item El sistema registra un evento de auditoría anónimo de la
    operación (CU-11: Registrar Evento de Auditoría).
\end{enumerate} \\
\textbf{Extensiones} &
\begin{itemize}[leftmargin=*, itemsep=2pt]
    \item[1a.] El personal prefiere dictar la respuesta por voz en
    lugar de escribirla: continúa en CU-05 (Dictar Respuesta por Voz),
    que extiende este caso de uso en el punto de extensión ``ingreso de
    texto''.
    \item[5a.] El paciente o el personal marcan la respuesta como
    incorrecta o ambigua: continúa en CU-10 (Reportar Traducción
    Incorrecta), que extiende este caso de uso en el punto de extensión
    ``revisión de mensaje''.
\end{itemize} \\
\textbf{Requisitos especiales} & RNF-T-05 (sincronización sin demora perceptible). \\
\textbf{Variantes de tecnología y datos} &
Ingreso de texto manual; sincronización entre estaciones mediante
BroadcastChannel (canal \texttt{nexora-chat}, nombre interno heredado del
repositorio de código). \\
\textbf{Frecuencia} &
Cada vez que el personal necesita responder al paciente durante la
atención. \\
\textbf{Temas abiertos} &
Sin historial persistente entre sesiones: si se cierra la pestaña, se
pierde la conversación. \\
\bottomrule
\end{longtable}

% ============================================================================
\subsection{CU-05: Dictar Respuesta por Voz}

\begin{longtable}{@{}>{\raggedright\arraybackslash}p{4.1cm} >{\raggedright\arraybackslash}p{\dimexpr\textwidth-4.1cm-2\tabcolsep\relax}@{}}
\caption{CU-05 --- Dictar Respuesta por Voz}\\
\toprule
\textbf{Nombre} & Dictar Respuesta por Voz \\
\textbf{ID} & CU-05 \\
\textbf{Actor principal} & Personal de Recepción o Enfermería \\
\textbf{Personal involucrado e intereses} &
Personal de Recepción o Enfermería: quiere responder sin escribir, por
ejemplo mientras atiende otras tareas. \\
\textbf{Precondiciones} &
El navegador soporta reconocimiento de voz; el personal está en el
punto de extensión ``ingreso de texto'' de CU-04. \\
\textbf{Garantías de éxito} &
El texto dictado queda disponible como respuesta, lista para enviarse
(el flujo continúa en CU-04 desde el paso de confirmación). \\
\textbf{Escenario principal de éxito} &
\begin{enumerate}[leftmargin=*, itemsep=2pt]
    \item El personal activa el dictado por voz.
    \item El sistema escucha el habla del personal de forma
    continua.
    \item El sistema transcribe el habla a texto en español rioplatense,
    mostrando resultados parciales a medida que reconoce.
    \item El personal detiene el dictado cuando termina de hablar.
    \item El sistema conserva el texto transcripto, sumándolo a cualquier
    texto que el personal ya hubiera escrito antes de dictar.
    \item El flujo continúa en el paso 2 de CU-04 (confirmar y enviar).
\end{enumerate} \\
\textbf{Extensiones} &
\begin{itemize}[leftmargin=*, itemsep=2pt]
    \item[1a.] El navegador no soporta reconocimiento de voz: el sistema
    informa la limitación y el personal continúa con el flujo
    principal de CU-04 (escribir).
    \item[2a.] Ocurre un error de micrófono o de reconocimiento: el
    sistema informa el error y detiene la escucha.
\end{itemize} \\
\textbf{Requisitos especiales} & RF-02.2; restricción de compatibilidad de navegador (ver ERS, secciones 3.4 y 3.6). \\
\textbf{Variantes de tecnología y datos} &
Web Speech API nativa del navegador, configurada en español argentino
(\texttt{es-AR}). \\
\textbf{Frecuencia} & Opcional, a criterio del personal de recepción o enfermería, en cualquier respuesta. \\
\textbf{Temas abiertos} &
Requiere Google Chrome; sin soporte confirmado en otros navegadores. \\
\bottomrule
\end{longtable}

% ============================================================================
\subsection{CU-06: Capturar Muestra de Seña para Entrenamiento}

\begin{longtable}{@{}>{\raggedright\arraybackslash}p{4.1cm} >{\raggedright\arraybackslash}p{\dimexpr\textwidth-4.1cm-2\tabcolsep\relax}@{}}
\caption{CU-06 --- Capturar Muestra de Seña para Entrenamiento}\\
\toprule
\textbf{Nombre} & Capturar Muestra de Seña para Entrenamiento \\
\textbf{ID} & CU-06 \\
\textbf{Actor principal} & Desarrollador de IA / Equipo de Proyecto \\
\textbf{Personal involucrado e intereses} &
Desarrollador de IA / Equipo de Proyecto: quiere ampliar el vocabulario
reconocible de forma simple. \newline
Paciente: se beneficia indirectamente de un vocabulario más amplio en el
futuro. \\
\textbf{Precondiciones} &
Estación del Paciente abierta en Modo Captura; cámara autorizada; backend
en ejecución. \\
\textbf{Garantías de éxito} &
Una o más muestras quedaron grabadas y asociadas a la etiqueta de seña
correspondiente. \\
\textbf{Escenario principal de éxito} &
\begin{enumerate}[leftmargin=*, itemsep=2pt]
    \item El desarrollador entra al Modo Captura.
    \item El desarrollador selecciona o crea la etiqueta de la seña a
    grabar (CU-07: Gestionar Vocabulario de Señas).
    \item El desarrollador indica cuántas repeticiones va a grabar de
    corrido.
    \item El desarrollador inicia la secuencia de grabación.
    \item El sistema anuncia una cuenta regresiva antes de cada
    repetición.
    \item El sistema graba los movimientos del desarrollador durante una
    ventana de tiempo fija.
    \item El sistema envía la muestra grabada al backend y la guarda
    asociada a la etiqueta.
    \item El sistema repite los pasos 5 a 7 hasta completar las
    repeticiones indicadas.
    \item El sistema confirma que las muestras quedaron guardadas y
    actualiza el conteo de muestras por etiqueta.
\end{enumerate} \\
\textbf{Extensiones} &
\begin{itemize}[leftmargin=*, itemsep=2pt]
    \item[2a.] La etiqueta no existe todavía: el sistema la agrega al
    vocabulario disponible en el momento de guardar la primera muestra.
    \item[6a.] El desarrollador detiene la grabación antes de tiempo: el
    sistema cancela esa repetición y no la guarda.
\end{itemize} \\
\textbf{Requisitos especiales} & RF-T-03, RF-T-04. \\
\textbf{Variantes de tecnología y datos} &
Modo legado (MediaPipe Holistic: se guardan solo landmarks) o Modo Tasks
(MediaPipe Tasks: se guardan landmarks y 52 coeficientes de expresión
facial); persistencia como archivos \texttt{.npy} en árboles de carpetas
separados por modo (\texttt{dataset/custom} vs.\ \texttt{dataset/custom\_tasks}
+ \texttt{dataset/custom\_blendshapes}). \\
\textbf{Frecuencia} &
Bajo demanda, cuando se necesita ampliar o reforzar el vocabulario. \\
\textbf{Temas abiertos} &
Sin validación automática de la calidad de la seña grabada: la calidad
depende del criterio del desarrollador. \\
\bottomrule
\end{longtable}

% ============================================================================
\subsection{CU-07: Gestionar Vocabulario de Señas}

\begin{longtable}{@{}>{\raggedright\arraybackslash}p{4.1cm} >{\raggedright\arraybackslash}p{\dimexpr\textwidth-4.1cm-2\tabcolsep\relax}@{}}
\caption{CU-07 --- Gestionar Vocabulario de Señas}\\
\toprule
\textbf{Nombre} & Gestionar Vocabulario de Señas \\
\textbf{ID} & CU-07 \\
\textbf{Actor principal} & Desarrollador de IA / Equipo de Proyecto (caso de uso incluido desde CU-06) \\
\textbf{Personal involucrado e intereses} &
Desarrollador de IA / Equipo de Proyecto: necesita saber qué señas ya
existen y cuántas muestras tiene cada una, para decidir qué grabar. \\
\textbf{Precondiciones} & Backend en ejecución. \\
\textbf{Garantías de éxito} &
El desarrollador cuenta con la lista actualizada de etiquetas conocidas y
con la cantidad de muestras disponibles por etiqueta. \\
\textbf{Escenario principal de éxito} &
\begin{enumerate}[leftmargin=*, itemsep=2pt]
    \item El sistema consulta las clases que el modelo actualmente
    reconoce.
    \item El sistema consulta cuántas muestras existen por etiqueta en el
    conjunto de datos activo.
    \item El sistema combina ambas fuentes en una única lista de
    etiquetas conocidas.
    \item El desarrollador elige una etiqueta existente, o crea una
    nueva.
\end{enumerate} \\
\textbf{Extensiones} &
\begin{itemize}[leftmargin=*, itemsep=2pt]
    \item[4a.] El desarrollador crea una etiqueta nueva: el sistema la
    valida (solo letras minúsculas, números y guion bajo, hasta 32
    caracteres) antes de aceptarla.
\end{itemize} \\
\textbf{Requisitos especiales} & RF-T-01, RF-T-02. \\
\textbf{Variantes de tecnología y datos} &
Las estadísticas de muestras se calculan sobre el árbol de datos
correspondiente al modo activo (legado o Tasks). \\
\textbf{Frecuencia} &
Cada vez que se abre el Modo Captura, y cada vez que se agrega una
etiqueta. \\
\textbf{Temas abiertos} & Ninguno relevado. \\
\bottomrule
\end{longtable}

% ============================================================================
\subsection{CU-08: Entrenar el Modelo de Reconocimiento}

\begin{longtable}{@{}>{\raggedright\arraybackslash}p{4.1cm} >{\raggedright\arraybackslash}p{\dimexpr\textwidth-4.1cm-2\tabcolsep\relax}@{}}
\caption{CU-08 --- Entrenar el Modelo de Reconocimiento}\\
\toprule
\textbf{Nombre} & Entrenar el Modelo de Reconocimiento \\
\textbf{ID} & CU-08 \\
\textbf{Actor principal} & Desarrollador de IA / Equipo de Proyecto \\
\textbf{Personal involucrado e intereses} &
Desarrollador de IA / Equipo de Proyecto: quiere que el modelo aprenda las
señas nuevas capturadas. \newline
Paciente y Personal de Recepción o Enfermería: se benefician de un modelo
más preciso y con más vocabulario. \\
\textbf{Precondiciones} &
Existen muestras de entrenamiento capturadas (CU-06) para al menos una
clase de seña. \\
\textbf{Garantías de éxito} &
Se generó una nueva versión del modelo de reconocimiento, entrenada con
el conjunto de datos disponible, lista para usarse en las próximas
sesiones. \\
\textbf{Escenario principal de éxito} &
\begin{enumerate}[leftmargin=*, itemsep=2pt]
    \item El desarrollador inicia el proceso de entrenamiento fuera de
    línea.
    \item El sistema releva las clases de seña disponibles en el conjunto
    de datos.
    \item El sistema separa las muestras en un conjunto de entrenamiento
    y uno de validación.
    \item El sistema entrena el modelo de reconocimiento durante un
    número fijo de rondas de aprendizaje.
    \item El sistema evalúa la exactitud del modelo sobre el conjunto de
    validación en cada ronda.
    \item El sistema conserva la versión del modelo con mejor exactitud
    de validación.
    \item El sistema deja disponible el modelo entrenado y su lista de
    clases para que el backend lo cargue en el próximo inicio.
\end{enumerate} \\
\textbf{Extensiones} &
\begin{itemize}[leftmargin=*, itemsep=2pt]
    \item[2a.] No se encuentran clases ni muestras en el conjunto de
    datos: el sistema informa la situación y no entrena.
\end{itemize} \\
\textbf{Requisitos especiales} & RF-T-05. \\
\textbf{Variantes de tecnología y datos} &
Aumento de datos por recorte aleatorio, ruido gaussiano y escalado
durante el entrenamiento; separación 80/20 con semilla fija para
reproducibilidad. \\
\textbf{Frecuencia} &
Bajo demanda del desarrollador, típicamente después de acumular nuevas
muestras de varias sesiones de captura. \\
\textbf{Temas abiertos} &
Proceso manual, sin disparo automático ni interfaz web; requiere acceso a
la línea de comandos del servidor. \\
\bottomrule
\end{longtable}

% ============================================================================
\subsection{CU-09: Otorgar Consentimiento Explícito de Cámara}
\textit{Diseñado --- pendiente de implementación. Cubre el requisito oficial RF-03.}

\begin{longtable}{@{}>{\raggedright\arraybackslash}p{4.1cm} >{\raggedright\arraybackslash}p{\dimexpr\textwidth-4.1cm-2\tabcolsep\relax}@{}}
\caption{CU-09 --- Otorgar Consentimiento Explícito de Cámara}\\
\toprule
\textbf{Nombre} & Otorgar Consentimiento Explícito de Cámara \\
\textbf{ID} & CU-09 \\
\textbf{Actor principal} & Paciente \\
\textbf{Personal involucrado e intereses} &
Paciente: necesita entender, antes de que se lo grabe, para qué se usa la
cámara y poder negarse. \newline
Sanatorio (Comité de Ética): necesita que ningún paciente sea grabado sin
autorización explícita. \\
\textbf{Precondiciones} &
El paciente abrió la Estación del Paciente; la cámara todavía no fue
activada. \\
\textbf{Garantías de éxito} &
La cámara se activa únicamente si el paciente otorgó el consentimiento; si
lo rechaza, la cámara permanece inactiva y CU-01 no puede comenzar. \\
\textbf{Escenario principal de éxito} &
\begin{enumerate}[leftmargin=*, itemsep=2pt]
    \item El sistema muestra la pantalla de consentimiento antes de
    solicitar permiso de cámara al navegador (Pantalla P-05, diseño
    nuevo).
    \item El sistema explica, en lenguaje simple, qué captura la cámara,
    para qué se usa y que no se guarda video ni imágenes (RNF-03 oficial).
    \item El paciente otorga el consentimiento.
    \item El sistema registra el consentimiento (marca temporal, sin
    datos identificatorios) y activa la cámara.
    \item El flujo continúa en el paso 1 de CU-01.
\end{enumerate} \\
\textbf{Extensiones} &
\begin{itemize}[leftmargin=*, itemsep=2pt]
    \item[3a.] El paciente rechaza el consentimiento: el sistema no
    activa la cámara y muestra una alternativa (por ejemplo, solicitar
    asistencia del personal); CU-01 no puede iniciar.
    \item[3b.] El navegador deniega el permiso de cámara después de un
    consentimiento otorgado: el sistema informa el error sin activar el
    flujo de reconocimiento.
\end{itemize} \\
\textbf{Requisitos especiales} & RF-03.1; RNF-03.1 (sin almacenamiento de video). \\
\textbf{Variantes de tecnología y datos} &
Diseño nuevo, sin componente real todavía; se prevé una pantalla modal
previa a \texttt{navigator.mediaDevices.getUserMedia()}. \\
\textbf{Frecuencia} & Una vez por sesión de la Estación del Paciente. \\
\textbf{Temas abiertos} &
Persistencia del consentimiento entre recargas de la misma sesión;
formato exacto del texto legal a mostrar (pendiente de validación con el
Comité de Ética del Sanatorio). \\
\bottomrule
\end{longtable}

% ============================================================================
\subsection{CU-10: Reportar Traducción Incorrecta o Ambigua}
\textit{Diseñado --- pendiente de implementación. Cubre el requisito oficial RF-06. Extiende a CU-01 y CU-04.}

\begin{longtable}{@{}>{\raggedright\arraybackslash}p{4.1cm} >{\raggedright\arraybackslash}p{\dimexpr\textwidth-4.1cm-2\tabcolsep\relax}@{}}
\caption{CU-10 --- Reportar Traducción Incorrecta o Ambigua}\\
\toprule
\textbf{Nombre} & Reportar Traducción Incorrecta o Ambigua \\
\textbf{ID} & CU-10 \\
\textbf{Actor principal} & Paciente, Personal de Recepción o Enfermería \\
\textbf{Personal involucrado e intereses} &
Paciente y Personal: quieren señalar cuando un mensaje traducido no
representa lo que se quiso decir, para que el equipo mejore el modelo.
\newline
Desarrollador de IA / Equipo de Proyecto: necesita casos reales
identificados para priorizar el reentrenamiento (CU-08). \\
\textbf{Precondiciones} &
Existe al menos un mensaje en la conversación (generado por CU-01 o
CU-04). \\
\textbf{Garantías de éxito} &
El reporte queda registrado, asociado al mensaje señalado y con su motivo,
sin exponer datos identificatorios del paciente. \\
\textbf{Escenario principal de éxito} &
\begin{enumerate}[leftmargin=*, itemsep=2pt]
    \item El actor revisa un mensaje de la conversación (punto de
    extensión ``revisión de mensaje'' de CU-01 o CU-04).
    \item El actor marca el mensaje como incorrecto o ambiguo.
    \item El sistema solicita un motivo breve (por ejemplo: ``seña mal
    reconocida'', ``oración no tiene sentido'', ``palabra faltante'').
    \item El actor confirma el reporte.
    \item El sistema registra el reporte, asociado al mensaje señalado, y
    confirma al actor que el caso quedó registrado.
    \item El flujo vuelve al punto de extensión de origen (CU-01 o
    CU-04).
\end{enumerate} \\
\textbf{Extensiones} &
\begin{itemize}[leftmargin=*, itemsep=2pt]
    \item[4a.] El actor cancela antes de confirmar: el sistema descarta
    el reporte y no lo persiste.
\end{itemize} \\
\textbf{Requisitos especiales} & RF-06.1; RNF-03.1 (reporte sin datos identificatorios ni imágenes). \\
\textbf{Variantes de tecnología y datos} &
Diseño nuevo; se prevé un ícono de reporte sobre cada burbuja de la
conversación en P-02 y P-03 (ver Entregable 04, mockups). \\
\textbf{Frecuencia} &
Bajo demanda, cada vez que el paciente o el personal detectan una
traducción incorrecta. \\
\textbf{Temas abiertos} &
Mecanismo de priorización de los reportes acumulados para el reentrenamiento
(CU-08); no definido todavía. \\
\bottomrule
\end{longtable}

% ============================================================================
\subsection{CU-11: Registrar Evento de Auditoría}
\textit{Diseñado --- pendiente de implementación. Cubre el requisito oficial RF-04. Subfunción incluida por CU-01 y CU-04, sin actor humano.}

\begin{longtable}{@{}>{\raggedright\arraybackslash}p{4.1cm} >{\raggedright\arraybackslash}p{\dimexpr\textwidth-4.1cm-2\tabcolsep\relax}@{}}
\caption{CU-11 --- Registrar Evento de Auditoría}\\
\toprule
\textbf{Nombre} & Registrar Evento de Auditoría \\
\textbf{ID} & CU-11 \\
\textbf{Actor principal} & Ninguno --- subfunción interna del sistema (ver \texttt{IDS.md}, Sección~1, sobre la decisión de no modelar un actor de auditoría). \\
\textbf{Personal involucrado e intereses} &
Dirección del Sanatorio y Equipo de Proyecto UTN (interesados oficiales
dueños de RF-04): necesitan métricas operativas anónimas para evaluar el
impacto del sistema (E08 --- Informe de Impacto Operativo, EDT), sin que
eso implique acceso a datos identificatorios del paciente. \\
\textbf{Precondiciones} &
CU-01 o CU-04 se está ejecutando y alcanzó un punto de inclusión de este
caso de uso. \\
\textbf{Garantías de éxito} &
Un evento de auditoría anónimo quedó registrado, con tipo de operación,
marca temporal y resultado, sin datos identificatorios ni imágenes o
video. \\
\textbf{Escenario principal de éxito} &
\begin{enumerate}[leftmargin=*, itemsep=2pt]
    \item El sistema completa una operación relevante (predicción,
    traducción, síntesis de voz, respuesta del personal o captura de
    muestra).
    \item El sistema arma un evento de auditoría anónimo con el tipo de
    operación, la marca temporal y el resultado (éxito o error), sin
    identidad del paciente ni contenido de imagen o video.
    \item El sistema persiste el evento en el registro de auditoría.
\end{enumerate} \\
\textbf{Extensiones} &
\begin{itemize}[leftmargin=*, itemsep=2pt]
    \item[3a.] La operación original terminó en error: el evento se
    registra igual, con resultado ``error'', para que la métrica
    operativa lo refleje.
\end{itemize} \\
\textbf{Requisitos especiales} & RF-04.1; RNF-03.1 (anonimato, sin imágenes ni video). \\
\textbf{Variantes de tecnología y datos} &
Diseño nuevo; se prevé persistencia basada en archivos, consistente con el
resto del sistema (ver Entregable 09, Modelo de Datos). \\
\textbf{Frecuencia} & Una vez por cada operación de sistema relevante. \\
\textbf{Temas abiertos} &
Mecanismo de exportación o consulta de las métricas acumuladas (E08,
Informe de Impacto Operativo); no forma parte de la interfaz de LSA Salud
en esta fase. \\
\bottomrule
\end{longtable}
